\section{対角化・Jordan標準形}
まず, 対角化可能性について述べる。次に, 対角化の基本的な計算をし, 最後に対角化の威力がわかる問題をとりあげる。
\subsection{対角化可能性}
よく知られた対角化可能条件は3つある。
\besha
\begin{prop} (対角化可能判定法)\\
$A\in M_n(\mb{C})$に対して次は同値である。
\ben
\item $A\in M_n(\mb{C})$は対角化可能である。
\item $A$の最小多項式は重根を持たない。
\item $A$の重複度$\mu_{\alpha}$の固有値$\alpha$に対する固有空間$V_{\alpha}$について$\dim{V_{\alpha}} = \mu_{\alpha}$が成立する。
\item $A$には1次独立な$n$個の固有ベクトルが存在する。
\een
\end{prop}
\ensha
\begin{egprob} (対角化可能であるための十分条件)\\
$A\in M_n(\mb{C})$が次のいずれかの条件をみたすとき$A$は対角化可能であることを示せ。\footnote{\cite{linalg1}演習問題108pより。}
\ben
\item $A$が$n$個の相異なる固有値を持つ。
\item $A^{m} = I_n$を満たす自然数$m$が存在する。
\item $A^2 = A$を満たす(冪等行列)。
\item
\item
\een
\end{egprob}%%～～～が対角化可能であることを示せ系の問題募集中！

\begin{egans}
(1):固有多項式が重根を持たないので最小多項式もそうであるから。(2):最小多項式は$x^{m}-1$を割り切ることから重根を持っていないことになるから。(3):最小多項式は$x(x-1)$を割り切るから(2)と同様。
\end{egans}


\newpage
\subsubsection{実対称行列の直交対角化可能性}
\tbox{
\begin{theo}
$A$が$n$次\tf{実}正方行列とする. $A$が直交対角化可能であるための\tf{必要十分条件}は, $A$が\tf{実対称行列であること}である.
\end{theo}
}{}
\prf まず$A$が直交対角化可能であるとする. すなわち, 直交行列$T$が存在して$T^{-1}AT$が対角行列$D$であるとする. このとき$A = TDT^{-1}$の転置を考えると
\[{}^t A = {}^t (T^{-1}) {}^t D {}^t T = ({}^t T)^{-1} D {}^t T = (T^{-1})^{-1} D t^{-1} =   TDT^{-1} = A\]
だから$A$は実対称行列である.\par
逆に$A$が実対称行列であるとする. $n$に関する帰納法で示す. $n=1$のときは自明. $n>1$とする.  $A$の固有値$\alpha$(\tf{注:}$A$は実対称なので$\alpha$は必ず実数)と, その実固有空間$V(\alpha)\subset \mb{R}^n$を取り, \aka{$V(\alpha)$の直交補空間$V(\alpha)^{\bot}$を取る.} $V(\alpha)$の正規直交基底$\set{v_1,\dots, v_m}$と$V(\alpha)^{\bot}$の正規直交基底$\set{w_1,\dots, w_{n-m}}$を取ると, この二つの基底の和は$\mb{R}^n$の正規直交基底となる. よって, $T := [v_1 v_2 \, \dots \, v_m\, w_1\, \dots \, w_{n-m}]$とすると$T$は直交行列である. $\mb{R}^n$の標準基底を$\set{e_i}$とするとき,
\bealn{
T^{-1}A T &= [T^{-1}Av_1\, \dots\, T^{-1}Av_m \, T^{-1}AT e_{m+1}\,\dots\, T^{-1}AT e_{n}] \\
&= [\alpha T^{-1}v_1\, \dots\, \alpha T^{-1} v_m\, T^{-1}AT e_{m+1}\,\dots\, T^{-1}AT e_{n}]\\
&= [\alpha e_1 \, \dots\, \alpha e_m\, T^{-1}AT e_{m+1}\, \dots\, T^{-1}AT e_{n}]\\
&=: \alpha E_m \oplus A_1
}
となる. $A$は対称行列で$T$は直交行列だから, 証明の最初の議論と同様に$T^{-1}AT$もまた対称行列である. よって$A_1$は実対称行列であるから, $A_1$が直交対角化できればよい. $A_1$のサイズは$n-m$なので, 帰納法の仮定で, ある$n-m$次直交行列$T_1$により$T_1^{-1}A_1T_1$は対角行列となる. よって, $P = T(E_m \oplus T_1)$とおくと\tf{$P$もまた直交行列}であって, $P^{-1} = (E_m \oplus T_1^{-1})T^{-1}$であり,
\bealn{
P^{-1}AP &= (E_m\oplus T_1^{-1}) T^{-1} A T (E_m \oplus T_1) \\
&= (E_m\oplus T_1^{-1}) (\alpha E_m \oplus A_1) (E_m \oplus T_1) \\
&= \alpha E_m\oplus T_1^{-1}A_1T_1
}
となるので, $A$も直交対角化可能である. \qed
\tbox{
\begin{theo}
$A$が$n$次実正方行列とする. $A$が直交対角化可能であるための必要十分条件ｈ, $A$の固有ベクトルのみからなる$\mb{R}^n$の正規直交基底が存在することである.
\end{theo}
}{}
\prf $A$が直交対角化可能とする. $P^{-1}AP = \mr{diag}(\alpha_1,\dots, \alpha_n)$が対角行列で$P = [p_1\, p_2\, \dots\, p_n]$が直交行列なら, \bolm{Ap_i = \alpha_i p_i}である. よって$p_i$はすべて$A$の固有ベクトルで, $P$が直交行列なので$\set{p_i}_{1\leq i\leq n}$は正規直交基底となり, 固有ベクトルからなる正規直交基底の存在が示せた. 逆は, この議論を逆に辿って$P$を構成すればよい. \qed

\begin{egprob}
$A$が実対称行列とする. $A$の異なる固有値$\alpha,\beta$に属する固有ベクトルを$v,w$とすると$v$と$w$は直交することを示せ.
\end{egprob}
\begin{egans} $\alpha\neq 0$とする. $\alpha \lan \alpha v,\beta w\ran  = \lan \alpha v, w\ran = $
\end{egans}





\subsubsection{Jordan標準形}


Jordan標準形の「固有値$\lambda_i$の部分」の中のJordan Blockの最大のサイズが$l_i$だとすると, \tf{最小多項式の次数は$\sum l_i$になる。}なぜなら, 固有値$\lambda_i$の部分は
\[
\begin{pmatrix}
\lambda_i &1 &0 &\cdots &0\\
0 & \lambda_i & 0 & \cdots  & 0\\
0 & 0 & \lambda_i & \cdots & 0\\
0 & 0& 0 & \cdots & 0
\end{pmatrix}
\]
という形をしていて(対角成分のひとつ上は1か0), $\mu_i$を$\lambda_i$の重複度とするとそこから$I_{\mu_i}$を差し引いて$l_i$乗したらJordan Blockがすべて消失するからだ。この行列論的な現象は大事であるから注意しておく:
\begin{prop}
$A$をサイズ$n$上三角行列で対角成分がすべて0であるものとするとき, $A^{n} = O$である.
\end{prop}

ブロック型の行列においては, 次のような状況が成り立つというのは予想されるが, しっかり示しておこう.
\begin{prop}
$A$を複素行列で$A=A_1\oplus A_2\oplus \dots \oplus A_n$というブロック型の行列とする. このとき, $A$が対角化可能であることと, すべての$A_k$が対角化可能であることは同値.
\end{prop}
\prf 同じ「型」を持つ二つの行列の積において, ブロックの外側にある成分は計算に\tf{全く影響しない}ということに注意せよ!(小さいサイズで適当に調べて見よ). そのような型同士の積は, 対応するブロック同士で積を取れば良い. \par
したがって, $P_k^{-1}A_k P_k$が対角行列になるなら, $P=P_1\oplus \dots \oplus P_n$は$P^{-1}AP$を対角化することが分かる. \par
逆に$A$が対角化可能とせよ. $A_k$をJordan標準形にする$P_k$を取る. このとき先と同じ$P$を取って$P^{-1}AP$を考えるとこれはJordan標準形になるが, $A$が対角化可能なのでこれは対角行列である(!). よって$P_k^{-1}A_kP_k$ももともと対角行列である. \qed
\begin{rem}
$P^{-1}AP$が対角行列だとしても, $P$に余計な成分があちこちに入っていると$P$と$A$が同じ型を持たないということが起きるので, 「逆に」以降の議論は注意しなければならない. この議論は同時対角化でも通用する.
\end{rem}


次のような手ごわい問題も過去問に存在する。
\begin{egprob} (H11数学I)
$V=\mb{C}^n$とする。$A\in M_{n}(\mb{C})$とする。$V$の部分ベクトル空間$W$は, 任意の$w\in W$に対して$Aw\in W$が成り立つとき, $A$-不変であるという。
\ben
\item $A$が対角化可能であるための必要十分条件は, $V$の$A$不変部分空間はつねに$A$不変な補空間を持つことであることを示せ。
\item $A$が対角化可能でかつ$A$の各固有値の重複度が1であるための必要十分条件は,$V$の$A$不変な部分ベクトル空間はつねにただ一つの$A$不変な補空間を持つことであることを示せ。
\een
\end{egprob}












\subsection{標準形の計算}

\begin{egprob} (H28基礎科目I)
$a$を実数とする. $A = \bmat{a&1&2\\ 0&1&0 \\ -2 & 0 & 0}$について以下の問に答えよ.
\ben
\item $A$の固有値を求めよ.
\item $A$が対角化可能となる実数$a$をすべて求めよ.
\een
\end{egprob}

\begin{egprob}
次の行列のJordan標準形を求めよ.
\[
\bmat{
1 & a & 0 \\
0 & 1 & 0\\
0 & 0 & 2
}
\]

\end{egprob}

\begin{egprob}
次の対称行列を直交行列により対角化せよ.
\[
\bmat{1 & 2\\ 2 & 1}
\]
\end{egprob}



\begin{egprob} (東工大H30基礎)
変数$x,y$に関する3次斉次多項式全体のなす$\mb{R}$上のベクトル空間を$V$とする. すなわち,
\[V = \set{ax^3 + bx^2y + cxy^2 + dy^3\mid a,b,c,d\in \mb{R}}\]
と定める.
\ben
\item $T$は線形変換であることを示せ.
\item $T$の固有ベクトルからなる$V$の基底を求めよ.
\een
\end{egprob}

\begin{egans}
(1)は読者に託す. (2)も対角化の要領でやればよいだろう. まず$V$の基底は$(x^3,x^2y,xy^2,y^3)$の順に取る.
\bealn{
T(x^3) &= 3x^3 + 3x^2y  \\
T(x^2y) &= (x+y)(2xy + x^2) = x^3 + 3x^2y + 2xy^2\\
T(xy^2) &= 2x^2y + 3xy^2 + y^3\\
T(y^3) &= 3xy^2 + 3y^3
}
より表現行列は次の通り.
\[
A = \bmat{
3 & 1 & 0 & 0 \\
3 & 3 & 2 & 0 \\
0 & 2 & 3 & 3 \\
0 & 0 & 1 & 3
}
\]
$|A-xI|$を頑張って計算する. めんどくさい....
\end{egans}


\newpage
\subsection{($\heartsuit$)同時対角化}
「$AX = XA$」とか可換性の式が出てくることが院試では非常に多い！！！ 可換なだけかと思いきや, これには「同時対角化可能」という概念が絡んでいる. とくに, 与えられた$A$に対し$XA=AX$を満たす$X$が何次元分だけあるかという問題は頻出である. \par
まず同時対角化可能性の主張は以下の通りである. 証明を覚えてほしい.
\tbox{
\begin{theo}
$A,B$を対角化可能な複素正方$n$次行列とする. このとき,
\[A,Bが同時対角化可能 \iff AB=BA\]
\end{theo}
}{title=同時対角化可能性}
\prf (すむーずぷりん氏のツイートを大変参考にした)
$(\Rightarrow)$は易しい. $(\Leftarrow)$を示す. $A$の固有値を$\lambda_1,\dots,\lambda_s$とし, 固有空間を$V_A(\lambda_i)$で表す.  $AB=BA$から次が言えるということを覚えておこう.
\lefba{\begin{prac} $\lambda$を$\lambda_i$のいずれかとする. $v\in V_A(\lambda)$なら$Bv \in V_{A}(\lambda)$であることを示せ.
\end{prac}}
よって$B(V_A(\lambda)) \subset V_{A}(\lambda)$が言える. $A$は対角化可能なので $V_A(\lambda_i)$ ($i=1,\dots, s$)の基底$\set{v_{i1},\dots, v_{i\mu_i}}$が$\mb{R}^n$の基底をなす. $P$を$i=1,2,\dots, s$の順にこの基底を列ベクトルに並べた行列とすると, $P^{-1}AP$は(いつものように！)対角行列であり, 左上から$\lambda_1$が$\mu_1$個, ..., $\lambda_s$が$\mu_s$個並ぶ. \par
この$P$で$P^{-1}BP$を考えると $P^{-1}BP = P^{-1}\sumib{ Bv_{11} &\dots& Bv_{s\mu_s}}$ となるが, $B v_{ij} \in V_{A}(\lambda_i)$なので$Bv_{ij}$は基底$\set{v_{i1},\dots, v_{i\mu_i}}$で表示される. その表示を$P^{-1}$で引き戻せば標準基底に戻るので, $\set{v_{i1},\dots, v_{i\mu_i}}$たちで取った基底の順番に注意すると,  $P^{-1}BP$はサイズ$\mu_i$の行列$B_i$の直和になっている:
\[
P^{-1}BP = B_1\oplus \dots \oplus B_s =  \bmat{B_1 & \\ & B_2 \\ && \ddots \\ &&& B_s }
\]
各$B_i$をJordan標準形にする行列$Q_i \in \mr{GL}_{\mu_i}(\mb{C})$が存在. このとき$Q = Q_1\oplus \dots \oplus Q_s$とすると$Q^{-1} = Q_1^{-1}\oplus \dots \oplus Q_s^{-1}$で, ブロックごとに計算すれば
\[Q^{-1}(P^{-1}BP)Q = \bmat{Q_1^{-1}B_1Q_1 & \\ & Q_2^{-1}B_2Q_2 \\ && \ddots \\ &&& Q_s^{-1}B_sQ_s }\]
であり, これは$B$のJordan標準形. $B$は対角化可能としているので, 実際は対角行列.
一方で, $I_k$がサイズ$k$の単位行列のとき
\[Q^{-1}(P^{-1}A P) Q =  \bmat{Q_1^{-1}(\lambda_1I_{\mu_1})Q_1 & \\
& \ddots \\ && Q_s^{-1}(\lambda_s I_{\mu_s})Q_s }\]
で, 各ブロックはまた$\lambda_i I_{\mu_i}$であり, $PQ$により同時対角化される.\qed


\begin{egprob} (\tf{同時固有ベクトル})
$AB=BA$であるとき, $A,B$は共通の固有ベクトルを持つことを証明せよ.
\end{egprob}
\begin{egans}
\chatgpthint{$A$の固有空間は$B$で不変です。}
\end{egans}




\begin{egprob} (東北大H29共通1)\label{tohoku-H29-C1}
$Z\in M_{3}(\mb{R})$に対して
\[V(Z) = \set{X\in M_3(\mb{R})\mid ZX=XZ}\]
とおく。
\ben
\item $A = \bep 8 & -9 & -2\\ 6 & -7 & -2 \\ -6 & 9 & 4 \enp$のとき, $A$の固有値と各固有値に対する固有ベクトルを求めよ。
\item $V(A)$の次元を求めよ。
\item $B=\bep -1 & -2 & 1\\ 4 & 5 & -2 \\ 4 & 4 & -1 \enp$のとき, $V(B)$の次元を求めよ。
\een
\end{egprob}

\begin{egans} 本問は直接計算でもなんとかなってしまうが, さすがに同時対角化の証明でやったようなことを覚えておいた方が良い. \\
(1): 各自計算せよ. 固有値は2, 1で, 固有空間$V_1$, $V_2$は
\[V_1 = \lan \sumib{1\\1 \\ -1} \ran,  V_2 = \lan \sumib{3\\ 2 \\ 0}, \sumib{1 \\ 0 \\ 3} \ran \]
(2): $X\in V(A)$のとき$AX=XA$である. $P = \sumib{1 & 3 & 1\\ 1 & 2 & 0 \\ -1 & 0 & 3}$により$P^{-1}AP = \mr{diag}(1,2,2)$である. 定理の証明で述べられたように, $v_1\in V_1$に対して$Xv_1\in V_1$, $Xv_2\in V_2$であるから $P^{-1}XP = \sumib{a&0&0\\ 0&b&c\\ 0&d&e}$という形になる. 逆にこの右辺の形をした行列$Y$を任意に取ると, その形から, これは$P^{-1}AP$と可換である. よって$PYP^{-1}$は$A$と可換である. 以上より,  $(k,l)$成分のみが1で他が0の行列を$E_{k,l}\in \mr{M}_{3}(\mb{R})$とすると
\[V(A) = P\lan E_{1,1}, E_{2,2}, E_{2,3}, E_{3,2}, E_{3,3} \ran P^{-1}\]
であるから5次元である. \\
(3): 真似してみよ.
\end{egans}
\begin{prac}
上の(1),(3)を解け.
\end{prac}

\begin{egprob} (東京大H29 A 第5問)
$A\in M_{4}(\mb{R})$, $A^2 - 2A + 2I_4 =0$, $V_{A} = \{ X\in M_4(\mb{R})\mid AX=XA \}$ このとき$\dim V_{A}$を求めよ。
\end{egprob}
\begin{egans}
$x^2 - 2x + 2 = 0$は実数解を持たないので$A$は単位行列の定数倍とはならない。よって$A$, および$A$に共役な行列の最小多項式は$x^2 - 2x + 2$である。これは$1\pm i$を根に持つから, $A$の固有値は$\lambda = 1 + i$, $\ovl{\lambda}$ですべてである。よって固有多項式$\det{(xI - A)}$としてあり得るのは
\[(x- \lambda)(x-\ovl{\lambda})^{3}, (x-\lambda)^2(x-\ovl{\lambda})^2, (x-\lambda)^{3}(x-\ovl{\lambda})\]
のいずれかだが, 左と右の可能性はない。なぜなら, 明らかに$\det{(xI - A)}\in \mb{R}[x]$だから。よって$\det{(xI - A)} = (x^2 - 2x + 2)^2$であり, 最小多項式は重根を持たないので$A$は対角化でき, 固有ベクトルを並べた行列$P$を取り
\[
P^{-1}AP =
\begin{pmatrix}
\lambda & 0 & 0 & 0\\
0& \lambda & 0 & 0\\
0 & 0 & \ovl{\lambda} & 0 \\
0 & 0 & 0& \ovl{\lambda}
\end{pmatrix}
\]
とできる. さて, $X \in V_A$なら
\[P^{-1}X P = \bmat{a& b \\  c&d \\ && e&f\\ && g& h},\quad (a,\dots, h\in \mb{C})\]
の形に表せる. 逆に右辺の形の行列$Y$を与えるとき, それは$P^{-1}AP$と明らかに可換であり, $A$は$PYP^{-1}$と可換である. よって, 複素の行列$X$で$AX=XA$を満たすものは複素$8$次元(実16次元)だけあると分かる. \par
ここからいつ実に収まるのか考えたくなるが,それはいろいろと難しかった. 東大の佐久間氏の院試シケプリを参考にしたところ, なんと\tf{複素化}をしていた. $\mb{R}$より$\mb{C}$のほうが嬉しいというものの, $\mb{R}$の議論を$\mb{C}$に変えるために複素化を行うというのは面白く, テンソルの威力を感じた\footnote{にしても, これがA問題であることは驚きである. }.  以下, 同氏のそのアイデアを述べる. \par
$\dim_{\mb{R}}V(A) = \dim_{\mb{C}} V(A)\tens_{\mb{R}}\mb{C}$であるから$V(A) \tens_{\mb{R}}\mb{C}$を調べればよい. $V(A)\tens_{\mb{R}}\mb{C}$はいかにも$\mr{M}_{4}(\mb{C})$における$\set{XA=AX}$を満たす$X$の全体の$\mb{C}$ベクトル空間$W$だと思われる. 実際それが正しい. なぜなら, $f:\mr{M}_{4}(\mb{R})\to \mr{M}_{4}(\mb{R})$を $f(X)= AX-XA$とするとき, $\ker{f} = V(A)$であることに注意して, 同一視$\mr{M}_{4}(\mb{R})\tens_{\mb{R}}\mb{C} \cong \mr{M}_{4}(\mb{C})$; $R \tens c \leftrightarrow cR$のもとで
\[W = \ker{(f\tens \mr{id}_{\mb{C}})} = (V(A)\tens \mb{C} )\oplus \mr{M}_{4}(\mb{R})\tens \set{0} = V(A)\tens \mb{C} \]
であるから. よって実8次元である. \qed



\begin{ans} (計算がしんどすぎてダメな方法)\\
$v_1\neq 0$を取り, $Av_1 = v_2$とし, $v_1,v_2,v_3$が一次独立になるように$v_3$を取ってから$Av_3 = v_4$とする.
\lefba{
\begin{claim}
$v_1,\dots, v_4$は$\mb{R}$上一次独立である.
\end{claim}
\prf これらで張られる空間が3次元であることは分かっている.
$c_1v_1 + c_2 v_2 + c_3 v_3 + c_4 v_4 = 0$を非自明な関係式とする. $c_4\neq 0$が必要である.  $A$を作用させると, $Av_2 = A^2v_1 = 2v_2 - 2v_1$, $Av_4 = 2v_4 - 2v_3$だから
\[-2c_2v_1 + (c_1+2c_2)v_2 - 2c_4v_3 + (c_3 + 2c_4)v_4 = 0 \]
もし$c_3 + 2c_4 = 0$なら$c_1=\dots = c_4 = 0$が分かる. そうではないとき, 上の式に$\dfrac{c_4}{c_3  +  2c_4}$をかけたものと最初の式は$c_4v_4$が共通するので, それを引くことで$v_1,v_2,v_3$の線形関係式を得るから, それらの係数は0でなければならない. したがって,
\[c_1 = \dfrac{-2c_2c_4}{c_3 + 2c_4},  c_2 = \dfrac{c_4(c_1 + 2c_2)}{c_3 + 2c_4}, c_3= \dfrac{-2c_4^2}{c_3 + 2c_4}\]
$a_i = c_i/c_4$とすると
\[a_1 = \dfrac{-2a_2}{a_3 + 2}, a_2 = \dfrac{a_1 + 2a_2}{a_3 + 2}, a_3 = \dfrac{-2}{a_3 + 2}\]
最後の式は$a_3^2 + 2a_3 + 2 = 0$で, 実数で成り立つ式ではない. よって矛盾.
\qed
}
したがって$P = \sumib{v_1 & v_2  & v_3 & v_4}$は可逆. $A^2 = 2A - 2I_4$だから
\[P^{-1}AP = \bmat{ 0 &-2 & & \\ 1&2&& \\ &&0&-2 \\ &&1&2
}\]
$V(A)\cong V(P^{-1}AP)$であるから上の行列$B:=P^{-1}AP$で同じ問題を調べればよい. $X=(x_{ij})_{i,j}\in\mr{M}_{4}(\mb{R})$とすると
\[XB = \bmat{ x_{12} & -2x_{11} + 2x_{12} & x_{14} & -2x_{13} + 2x_{14} \\
x_{22}& -2x_{21} + 2x_{22}& x_{24} & -2x_{23} + 2x_{24}\\
x_{32} & -2x_{31} + 2x_{32} & x_{34} & -2x_{33} + 2x_{34} \\
x_{42}& -2x_{41} + 2x_{42}& x_{44} & -2x_{43} + 2x_{44}
}\]
\[BX = \bmat{ -2x_{21} & - 2x_{22} &-2x_{23}& - 2x_{24}\\
x_{11} + 2x_{21}& x_{12} + 2x_{22}& x_{13} + 2x_{23}& x_{14} + 2x_{24} \\
-2x_{41} & - 2x_{42} &-2x_{43}& - 2x_{44}\\
x_{31} + 2x_{41}& x_{32} + 2x_{42}& x_{33} + 2x_{43}& x_{34}+ 2x_{44} \\


}\]
ここから次元を計算すればできる. でもできない...
\end{ans}

\begin{comment}
$\Lambda = P^{-1}A P $であるとすると, $X\in M_{4}(\mb{R})$に関して
\[XA = AX \iff (P^{-1}XP)\Lambda = \lambda (P^{-1}XP)\]
だから
\beqn
V_A &=& \set{ X \in M_4(\mb{R}) \mid (P^{-1}XP)\Lambda = \Lambda (P^{-1}XP)}\\
&=& \set{X \in M_4(\mb{C}) \mid (P^{-1}XP)\Lambda = \Lambda (P^{-1}XP), X = \ovl{X}}
\eeqn
$Y\in M_{4}(\mb{C})$が$Y\Lambda = \Lambda Y$を満たすとして
\[
Y =
\begin{pmatrix}
A_{11} & A_{12} \\
A_{21} & A_{22}
\end{pmatrix}
\]
であるとすると
\[\lambda A_{12} = \ovl{\lambda}A_{12}, \lambda A_{21} = \ovl{\lambda}A_{21}\]
だから$A_{12} = A_{21} = O$であり,
\[
X=
P
\begin{pmatrix}
A_{11} & O \\
O & A_{22}
\end{pmatrix}
P^{-1}
\]
を得る。
\end{comment}

\end{egans}


\begin{egprob} (仮想面接問題) \\
対角化可能な$A$に対し, $\set{XA = AX}$なる$A$と同サイズの$X$の全体の$\mb{C}$ベクトル空間の次元は固有値の重複度を用いて表すとどのようになりますか?
\end{egprob}






\subsection{演習問題Part2}
\subsubsection{Level A}
\begin{prob} (2019前期 微分方程式レポート)\\
$A=
\bep
6 & -3 & -7\\
-1 & 2 & 1\\
5 & -3 & -6
\enp$
を対角化し, $\ve{x}(t) = \bep x_1(t)\\ x_2(t)\\ x_3(t) \enp$に対する微分方程式$\dfrac{d\ve{x}}{dt} = A\ve{x}$を解け。
\end{prob}


\begin{prob} (H28基礎科目II,2)
$A,B$を複素三次正方行列とする. $A$の最小多項式は$x^3-1$, $B$の最小多項式は$(x-1)^3$とする. このとき$AB\neq BA$であることを示せ.
\end{prob}

%=======================%
\subsubsection{Level B}

\begin{prob} (H21 RIMS 基礎)
$A\in \mr{M}_{n}(\mb{C})$に対し, $X\in \mr{M}_n(\mb{C})$で$AX=XA$なるもの全体の複素ベクトル空間$V$の次元は$n$以上であることを証明せよ.
\end{prob}


\begin{prob} (東北大H28共同) \href{https://drive.google.com/file/d/1anaTAHLXhDH8Ng_6pvw6pPFb4MDD5v7i/view?usp=sharing}{解答}
$3$次実正方行列
\[
A = \bep 0 & 2 & 1\\ 2 & 3 & 2 \\ 1 & 2 & 0 \enp
\]
に対して, 以下の問いに答えよ。
\ben
\item $A$の固有値をすべて求めよ。
\item (1)で求めた固有値に対し, その固有値に対する固有空間の基底を一組求めよ。
\item $n$を正の整数とするとき, $A^{n} \bep 1\\ 0 \\ 0 \enp$を$n$で表せ。
\item 次の2条件を満たす$3$次実正方行列$B$が存在することを示せ。
\[AB=BA,\quad B^2 = A\]
\een
\end{prob}

\begin{prob} (H30東北大共同1)\label{prob:tohoku_H30_kyodo1}
$a$を実数, $A$を行列
\[\bep
2 & a & -1\\
0 & 2 & 1\\
-1 & 8 & -1
\enp\]
に対して, 線形写像$f_{A}:\mb{R}^{3}\to\mb{R}^{3}$を$f_A(v)=Av$で定める。
\ben
\item 行列$A$の固有多項式$p_{A}(x)$を求めよ。
\item 行列$A$を複素行列によって対角化できないような実数$a$をすべて求めよ。
\item $a$を(2)で求めたものとする。$\mb{R}^{3}$の一次元部分空間$W$で
\ben[label=(\alph*)]
\item $f_{A}(W)  \subset  W$
\item $f_A$が誘導する線形写像$\witi{f}_{A}:\mb{R}^{3}/W\to \mb{R}^{3}/W$の表現行列は対角化可能である。
\een
\een
\end{prob}



\subsubsection{Level C}
\begin{prob} (H31基礎科目5) \label{prob:H31kiso5}
$n$を2以上の整数, $A$を$n$次複素正方行列とする. $A^{n-1}$は対角化可能でないが$A^n$が対角化可能であるとき, $A^n=O$であることを示せ.
\end{prob}



\subsection{Hint・Answer}
{\small
\begin{probans}
$x' = kx$は$x=ce^{kx}$が解. 対角化してから微分方程式を解き, 基底を戻す.
\end{probans}

\begin{probans}
$A$は対角化可能. もし$AB=BA$なら$B$も対角化可能で, 最小多項式に関する条件に反する.\\
\fbox{同時対角化の知識を避けたい場合} $AB=BA$とする. $\omega$が1の3乗根, $v_k$を固有値$\omega^k $に属する$A$の固有ベクトルとする. このとき$A(Bv_k) = \omega^k Bv_k$なので$Bv_k$も$A$の$\omega^k$に属する固有ベクトルとなる. $V_A(\omega^k)$は一次元なので$Bv_k = b_kv_k$なる$b_k \in \mb{C}$が存在. $b_k$は$B$の固有値なので$b_k=1$である. よって$v_0,v_1,v_2\in V_B(1)$. $\set{v_0,v_1,v_2}$は$A$では異なる固有値に属する固有ベクトルなので, 一次独立である. よって$P=[v_0 \, v_1\, v_2]$は正則で, $P^{-1}BP = P^{-1}[Bv_0\, Bv_1\, Bv_2] = P^{-1}[v_0\, v_1\, v_2] = I$なので 最小多項式は$x-1$でなければならない. これは矛盾. (\tf{ポイント:} $A$を対角化する行列$P$で$P^{-1}BP$を考えること.)
\end{probans}

%==================================Level B================

\begin{probans}
$A$がJordan標準形の場合に示せばよくなる. なぜか. さらに, $A$がJordan Blockの形の行列にして示せば一般の場合にも示せる. なぜか. さらに, $A=\lambda I + A'$と分ければ, $A'X = XA'$となるから$\lambda = 0$の場合のみ見ればよい. この特別な$A$に対して主張が示せるか?
\end{probans}



\begin{probans}
\chatgpthint{$A$を直交対角化して各固有空間上で平方根を選んでください。}
\end{probans}

\begin{probans}
\chatgpthint{固有多項式の重根と固有空間の次元を比較してください。}
\end{probans}

\begin{probans}
\chatgpthint{各Jordanブロックをべき乗したときの第1超対角成分を見てください。}
\end{probans}


%===========Level C===================

\begin{probans}
$A$をJordan標準形に取り替えてよい.\par  $J(c,m)$を\tf{サイズ2以上の}Jordan Blockの任意の一つとする.  $A^n$は対角化可能なので, そのBlockとしてある$J(c,m)^n$も対角化可能. よって$J(c,m)^n$の最小多項式に重根がないが, $J(c,m)^n$の固有値は$c^n$のみだから最小多項式は$T-c^n$である. よって$J(c,m)^n = c^n I$を得るが, 1行2列の成分比較から$nc^{n-1}=0$なので$c=0$である. よってBlockを並び替えると(並び替える操作は共役類を保つ), $A=D \oplus N$あるいは$A=N$という形にしてよい. ここで$D$は正則な対角行列で, $N$は$J(0,m)$の形の行列のいくつかの直和であるが, $A$にノンゼロな固有値が元から無い場合には$A=N$である. さて, 仮に$D$のサイズが1以上であったなら$N$のサイズは$n-1$以下だから, $N^{n-1} = O$. よって$A^{n-1} = D^{n-1}\oplus N^{n-1} = D^{n-1}\oplus O$は対角行列なので仮定に矛盾. よって$A=N$だが確かに$A^n = O$である. \qed \\

\begin{comment}
 $A$の固有値を$\lambda_i$ ($i=1,\dots, s$)とする. このとき, $s=1$である. [証明] $s\geq 2$ならJordan Blockはすべてサイズ$n$未満である.  $A^{k}$は常にブロック型の行列である. $A^{n-1}$が対角化可能でないので, あるブロック$J(c,m)$が$n-1$乗で対角化可能でない. しかし$A^n$が対角化可能なので, そのブロック$J(c,m)$は$n$乗で対角化可能である. よって$J(c,m)^n = c^nI$だが, 左辺の1行2列目が$nc^{n-1}$なので$c=0$でなければならない. しかし$m<n$なので$J(c,m)^{n-1} = O$であり, 対角行列なのでおかしい.[終] \par
 よって$s=1$であり, $A$の固有値は1種類しかないので, それを$c$とおく.  このとき, $A^n = c^nI$である. $A$のJordan Blockが$r$個あるとする. このとき$r=1$である. [証明] $r\geq 2$とせよ. Block $J(c,m)$を任意に一つ取ると$m<n$である. $A^n = c^n I$により$J(c,m)^n = c^n I$であるべきだが, 先と同様の成分比較で$c=0$でなければならない. しかしこのとき$m<n$により$J(0,m)^{n-1} = O$である. $J(c,m)$は任意のJordan Blockなので$A^{n-1}$が対角化可能となり不適. [終] よって$r=1$なので$A=J(c,n)$と書ける. $A^{n} = c^n I$で同様の成分比較により$c=0$を得る. このとき確かに$A^n = J(0,n)^{n} = O$である.
 \end{comment}

\end{probans}
}




\clearpage
